\chapter{The complex}
\label{ch:complex}

\texttt{AMRComplex} builds a single cell complex over the whole plotfile
hierarchy. Everything in this chapter follows from one decision, stated first.

\section{Identity is an index, not a position}
\label{sec:identity}

A node is the pair $(\lev,\iv)$: the level and the integer index of the cell
whose centre it occupies. It is stored as a plain aggregate so that it is
trivially copyable and can therefore live inside a \texttt{BaseFab} and travel
through \texttt{FillBoundary} and \texttt{ParallelCopy}.

\begin{designrule}\label{rule:identity}
Two elements share a vertex because they name the same cell, never because two
coordinates compared close.
\end{designrule}

The ordering on nodes is lexicographic---by level, then by index---which is a
genuine strict weak ordering, unlike the $\epsilon$ comparison of
\cref{ch:assessment}. Three substantial consequences follow, and they are the
substance of the design rather than incidental benefits.

\section{Periodic identification is free}
\label{sec:periodic}

\texttt{FillBoundary} copies node \emph{values} into ghost cells. It cannot
shift an index stored inside the object it copies. Therefore, after a periodic
fill, the ghost beyond the high face literally contains the low-face node. That
copy \emph{is} the periodic identification: there is no arithmetic, no modular
reduction, and no opportunity for the coordinate-based defect of
\cref{ch:assessment}.

This satisfies~\ref{req:periodic} exactly. No boundary plane enters any
calculation, so the arbitrary placement of the domain faces cannot affect any
reported quantity---which is why the doubly and triply periodic cases cost
nothing here, a point that returns decisively in \cref{ch:topology}.

\section{Coarse--fine stitching}
\label{sec:stitch}

The construction is that of \cref{ch:assessment}: fine ghosts lying outside the
level are stamped with their coarse parent's node, so that a block of
$r^{\dsp}$ fine positions resolves to one coarse node and interface elements
become degenerate cells. Refinement ratio is unrestricted; ratios $2$ and $4$
are both verified.

\paragraph{Proper nesting is a precondition, not a nicety.}
A level's ghost ring must find valid parents one level down. If the parent level
is itself covered there by this level, the stitch has nothing to attach to and
the interface silently leaks volume. The first three-level test failed with a
$3\%$ residual for exactly this reason---a defect in the test generator, but one
that nothing detected. Construction now aborts when a ghost stamped from level
$\lev-1$ comes back covered, naming the levels involved.

\section{Element ownership}
\label{sec:ownership}

Elements must be generated exactly once, globally, without a deduplication
pass. The obvious rule is wrong, and wrong in a way that silently removes
exactly the elements the design exists to produce.

Consider a one-dimensional slice with refinement ratio $2$, a fine region
occupying coarse cells $[8,24)$ and therefore fine indices $[16,48)$. Coarse
cell $7$ is valid, at $7.5\,\dvol_c^{1/\dsp}$; coarse cell $8$ is covered; the
first fine cell centre is at $8.25$ in coarse units. The bridging element spans
$[7.5, 8.25]$ and its base cell is fine index $15$---a \emph{ghost}. No fine
box holds it as valid, and the coarse box that would hold it has a covered
corner and is rejected. Under the rule ``the box whose valid region contains the
base cell'', nobody generates it, leaving a gap one stitch-layer thick around
every refined region. Measured on a two-level $32^3$ case: \textbf{1.93\%\ of
the domain volume}.

The rule that works: the base range is the valid box grown by one \emph{on the
low side}, and a base cell in that shell belongs to the lowest-indexed box that
would visit it. Since a \texttt{BoxArray} is globally replicated this is
deterministic and requires no communication: the same element is produced by
exactly one rank at any rank count. Periodic wrap images are excluded, since the
identical element---the same node set, identity being topological---is already
produced where the cell is valid.

Ownership is resolved as box algebra once per grid, not as a query per cell: the
low-side shell is at most $\dsp$ boxes, the valid-elsewhere and periodic-image
exclusions are \texttt{BoxList} subtractions, and the lower-index claim is one
\texttt{intersections} call.

\section{Global numbering}
\label{sec:numbering}

A \texttt{BoxArray} is a globally replicated object: \texttt{ba[i]} is the same
box on every rank for any number of ranks, and only the
\texttt{DistributionMapping} depends on the rank count. Its own index order is
therefore already a deterministic global ordering, and
\begin{equation}
  \gid(\lev,\iv) \;=\; O_{\lev} \;+\; o_{\lev,i} \;+\; \mathrm{lex}(\iv, B_{\lev,i}),
  \label{eq:gid}
\end{equation}
where $O_{\lev}$ is the cumulative cell count of levels below $\lev$,
$o_{\lev,i}$ the cumulative count of boxes below $i$ within level $\lev$, and
$\mathrm{lex}$ the lexicographic offset within box $B_{\lev,i}$.

\Cref{eq:gid} needs \emph{no communication at all}: every term is a box size,
known globally. It counts every cell rather than only the uncovered ones, which
costs nothing---no array is allocated per identifier---and buys a great deal:
$\gid$ becomes a pure function of $(\lev,\iv)$, so it can be assigned at fill
time and then travels with the identity through \texttt{FillBoundary}, with no
second pass to propagate identifiers into ghosts. Identifiers are consequently
sparse over valid nodes; a compaction pass can be added for formats that need
them contiguous.

\section{Field data}
\label{sec:fields}

An arbitrary vector of fields is carried on the nodes, on the same grids with
the same ghost width as the identity, and propagated by the same calls: one
\texttt{FillBoundary} with the same periodicity, and the coarse-parent stamp
performing identity and value \emph{in the same loop under the same test}.

Stamping them in a second loop would be a second opportunity for a stitch
node's value to disagree with the identity it belongs to, and nothing
downstream would notice: the geometry would still verify and the extracted
surface would simply be wrong. Fields are initialised to a sentinel so that an
unpropagated value is detectable rather than plausible, and the verification
counts valid nodes still carrying it.

\section{The canonical traversals}

Two template traversals are the only sanctioned access:
\texttt{forEachElement}, which yields each element's corner nodes and their
field values, and \texttt{forEachNode}, which yields each valid node once with
its cell volume and a flag marking nodes another rank may also see. Every
consumer---verification, contouring, segmentation---goes through them, so that
ownership, corner ordering and the validity test are stated once. Two copies of
such a loop are two opportunities to disagree about which elements exist.
